% Generated draft: requires mathematical and source review.
\documentclass[11pt]{article}
\usepackage[margin=25mm]{geometry}
\usepackage{fontspec}
\setmainfont{Latin Modern Roman}
\usepackage{amsmath,amssymb,mathtools,unicode-math}
\setmathfont{Latin Modern Math}
\usepackage{enumitem,needspace,xurl,hyperref}
\setlength{\parindent}{0pt}
\setlength{\parskip}{5pt}
\setlength{\emergencystretch}{3em}
\newcommand{\N}{\mathbb{N}}
\newcommand{\Z}{\mathbb{Z}}
\newcommand{\Q}{\mathbb{Q}}
\newcommand{\R}{\mathbb{R}}
\newcommand{\C}{\mathbb{C}}
\newcommand{\F}{\mathbb{F}}
\newcommand{\A}{\mathbb{A}}
\newcommand{\PP}{\mathbb{P}}
\newcommand{\E}{\mathbb{E}}
\newcommand{\eps}{\varepsilon}
\newcommand{\dd}{\,\mathrm{d}}
\begin{document}
\section*{UnsolvedMath ID: 30003166}
\textbf{Source problem code:} catalogue entry
\subsection{Definitions and mathematical statement}
An \emph{arithmetic scheme} means a separated scheme of finite type over \(\operatorname{Spec}\Z\). A complex of abelian groups is \emph{perfect} if it is quasi-isomorphic to a bounded complex of finitely generated projective \(\Z\)-modules. A complex of real vector spaces is perfect if it is quasi-isomorphic to a bounded complex of finite-dimensional real vector spaces.

\textbf{Conjectural existence problem, with unresolved source parameters.} For every arithmetic scheme \(X\), the source asks for compactly supported Weil--\'{e}tale and Weil--Arakelov cohomology complexes
\[
C_{W,c}(X,n):=R\Gamma_{W,c}(X,\Z(n)),\qquad
C_{\mathrm{ar},c}(X,A,n):=R\Gamma_{\mathrm{ar},c}(X,A(n)).
\]
The supplied record does not specify the range or quantification of the twist \(n\), the allowed coefficients \(A\), or the full intended derived categories. These remain parameters requiring clarification; no additional restrictions on \(X\) are supplied.

The requested complexes must satisfy at least the following conditions:
\begin{enumerate}
\item \(C_{W,c}(X,n)\) is a perfect complex of abelian groups.
\item \(C_{\mathrm{ar},c}(X,\R,n)\) is a perfect complex of real vector spaces.
\item Cup product with the source's fundamental class
\[
\theta\in H^1(\operatorname{Spec}\Z,\R(0))
\]
makes the graded cohomology of \(C_{\mathrm{ar},c}(X,\R,n)\) into an acyclic complex. Explicitly, writing
\[
V^i=H^i\bigl(C_{\mathrm{ar},c}(X,\R,n)\bigr),
\]
the intended cup-product action must give maps \(d_\theta^i:V^i\to V^{i+1}\) such that, for every \(i\in\Z\),
\[
d_\theta^{i+1}\circ d_\theta^i=0,
\qquad
\ker(d_\theta^i)=\operatorname{im}(d_\theta^{i-1}).
\]
This is acyclicity of the cohomology groups equipped with the cup-product differential; it does not assert that each \(V^i\) vanishes.
\end{enumerate}

\textbf{Unresolved qualifications.} The cohomology theory containing \(\theta\), its normalization, and its action on compactly supported Weil--Arakelov cohomology are not defined in the supplied record. The displayed conditions alone do not characterize the intended cohomologies: zero complexes with zero action satisfy them. Additional structural requirements are therefore needed to distinguish the intended existence problem from a bare existence assertion about perfect complexes. The background mentions zeta-value properties, but provides no zeta-value formula or compatibility condition to state.
\subsection{Short English statement}
The supplied conjecture asks for compactly supported Weil--\'{e}tale and Weil--Arakelov cohomologies for every separated scheme of finite type over \(\Z\). The integral Weil--\'{e}tale complex and the real Weil--Arakelov complex should be perfect, and cup product with a fundamental class should give an exact complex on the latter's cohomology. The twist range, coefficient choices, and structural requirements identifying these cohomologies are missing. This is a statement-only draft requiring human mathematical review; the supplied open status is not independently verified.
\subsection{Sources}
\begin{itemize}
\item Supplied local UnsolvedMath registry record, permanent ID \(30003166\), source problem code \texttt{catalogue entry}, in the Oberwolfach Reports Open Problems collection. The record gives proposed year 2016 and names no individual proposer.
\item Supplied report citation: \emph{Algebraic K-theory and Motivic Cohomology (2017)}, Oberwolfach Reports, pp.\ 1753--1807, DOI \texttt{10.4171/owr/2016/31}. \url{https://doi.org/10.4171/owr/2016/31}. The collection describes Oberwolfach Reports as documenting workshops and their open-problem discussions. The supplied citation's 2017 title date is retained alongside its 2016 DOI and proposed-year metadata without reconciliation.
\item Flach--Morin, \emph{Weil--Etale Cohomology and Zeta-Values}. Supplied literature link: \url{https://ems.press/journals/dm/articles/8965580}.
\item Beshenov, \emph{Weil--etale cohomology at negative integers}. Supplied literature link: \url{https://arxiv.org/abs/2102.12114}.
\item Supplied literature assessment, checked exactly \texttt{2026-08-21}: the record labels the problem open and describes compactly supported Weil--Arakelov and Weil--\'{e}tale cohomology and the listed zeta-value properties as conjectural in the stated arithmetic-scheme generality, while proved or compatible in several restricted cases. This is supplied dated curation derived from the OWR material, not an independently verified assessment. The restricted cases and zeta-value properties are not specified in this record; time-sensitive status and cited claims require independent verification.
\item The registry labels its statement \texttt{corrected\_verified} and reports that it was checked against the official Oberwolfach report, preserving the report's hypotheses, notation, and open target. These are supplied verification claims; no source visit or independent verification is asserted here.
\end{itemize}
\end{document}
